\documentclass[a4paper]{article}
\usepackage[fontset=fandol]{ctex}
\usepackage{amsmath, amssymb}
\usepackage{graphicx}
\usepackage[margin=2.5cm]{geometry}
\usepackage[most]{tcolorbox}
\usepackage{etoolbox}
\usepackage{listings}
\usepackage{hyperref}
\usepackage{booktabs}
\usepackage{subcaption}
\usepackage{float}
\usepackage{tikz}
\IfFileExists{pgfplots.sty}{\usepackage{pgfplots}\pgfplotsset{compat=1.18}}{}
\newtcolorbox{knowledgebox}[1]{enhanced, colback=blue!5!white, colframe=blue!75!black, colbacktitle=blue!75!black, coltitle=white, fonttitle=\bfseries, title=#1, attach boxed title to top left={yshift=-2mm, xshift=2mm}, boxrule=1pt, sharp corners}
\newtcolorbox{importantbox}[1]{enhanced, colback=yellow!10!white, colframe=yellow!80!black, colbacktitle=yellow!80!black, coltitle=black, fonttitle=\bfseries, title=#1, sharp corners}
\newtcolorbox{warningbox}[1]{enhanced, colback=red!5!white, colframe=red!75!black, colbacktitle=red!75!black, coltitle=white, fonttitle=\bfseries, title=#1, sharp corners}
\lstset{language=Python,basicstyle=\ttfamily\small,keywordstyle=\color{blue},stringstyle=\color{red!60!black},commentstyle=\color{green!60!black},breaklines=true,frame=single,numbers=left,numberstyle=\tiny\color{gray},captionpos=b,extendedchars=false}
\newcommand{\notetitle}{Advanced Topics on End-to-End ASR I}
\newcommand{\noteauthors}{Codex 基于课程页、公开视频字幕、coverage ledger 与补充官方材料整理}
\newcommand{\notedate}{\today}
\newcommand{\videochannel}{WAVLab}
\newcommand{\videopublishdate}{2023-11-15}
\newcommand{\videoduration}{unknown}
\newcommand{\videourl}{https://www.youtube.com/watch?v=08BIoO2n-TM}
\newcommand{\videocoverpath}{cover.jpg}
\begin{document}
\begin{titlepage}
\centering
{\Large 课程笔记\par}
\vspace{1.2cm}
{\huge\bfseries \notetitle\par}
\vspace{0.8cm}
{\large \noteauthors\par}
\vspace{0.3cm}
{\large \notedate\par}
\vspace{1.2cm}
\ifdefempty{\videocoverpath}{{\small 公开视频封面不可用。\par}}{\includegraphics[width=0.82\textwidth,height=0.45\textheight,keepaspectratio]{\videocoverpath}\par}
\vfill
\begin{tcolorbox}[width=0.9\textwidth, colback=black!2!white, colframe=black!60, sharp corners]
\textbf{视频作者/频道}：\videochannel\par
\textbf{发布日期}：\videopublishdate\par
\textbf{视频时长}：\videoduration\par
\textbf{视频链接}：\href{\videourl}{\nolinkurl{\videourl}}
\end{tcolorbox}
\end{titlepage}
\tableofcontents
\newpage
\section{Data variation and orientation}
根据课程页主题、公开视频字幕以及本讲的 coverage ledger，Data variation and orientation 这一部分主要围绕 Advanced Topics on End-to-End ASR I | Advanced topics on end-to-end ASR I；speech data 的 variation 来源、data orientation，以及为什么鲁棒性是 advanced ASR 的起点。 展开。
写作上保持 coverage-first 原则：先把老师在这一阶段真正讲过的问题设定、模型部件、训练或推理逻辑讲清，再压缩成适合复习的结构。

\begin{figure}[H]
\centering
\includegraphics[width=0.92\textwidth]{figures/coverage_map.png}
\caption{本讲 coverage map：按 coverage\_units.jsonl 归纳出的核心主题与章节映射。}
\end{figure}

\begin{figure}[H]
\centering
\includegraphics[width=0.92\textwidth]{figures/timeline.png}
\caption{本讲 timeline：根据字幕时间范围归纳的主题推进顺序。}
\end{figure}

\subsection{Advanced Topics on End-to-End ASR I}
从证据位置看，这一部分对应的课程证据区间是 2023-11-15。课程在这里强调的核心内容可以概括为：Advanced Topics on End-to-End ASR I | Advanced topics on end-to-end ASR I。
从课程组织上看，这一部分承担的是把局部概念嵌回整条识别流水线中的作用，避免读者把公式或模块孤立地记忆。

在本章结构中，这一单元被并入“Data variation and orientation”，目的是把同一阶段的定义、机制和权衡集中讲清。

\subsection{speech data 的 variation 来源、data orientation，以及为什么鲁棒性是 advanced ASR 的起点}
从证据位置看，这一部分对应的课程证据区间是 00:00:00,120--00:07:05,000。课程在这里强调的核心内容可以概括为：speech data 的 variation 来源、data orientation，以及为什么鲁棒性是 advanced ASR 的起点。。
从课程组织上看，这一部分承担的是把局部概念嵌回整条识别流水线中的作用，避免读者把公式或模块孤立地记忆。

在本章结构中，这一单元被并入“Data variation and orientation”，目的是把同一阶段的定义、机制和权衡集中讲清。

\subsection{本章小结}
Data variation and orientation 这一部分把 Advanced Topics on End-to-End ASR I 中相邻的教学单元串成了一条连续叙事：既交代了这一阶段要解决的问题，也交代了方法、推理或工程权衡为什么要这样组织。

\section{Augmentation and robustness}
根据课程页主题、公开视频字幕以及本讲的 coverage ledger，Augmentation and robustness 这一部分主要围绕 SpecAugment / noise perturbation / augmentation 技术如何改善 ASR 训练。 展开。
写作上保持 coverage-first 原则：先把老师在这一阶段真正讲过的问题设定、模型部件、训练或推理逻辑讲清，再压缩成适合复习的结构。

\subsection{SpecAugment / noise perturbation / augmentation 技术如何改善 ASR 训练}
从证据位置看，这一部分对应的课程证据区间是 00:07:02,400--00:13:38,560。课程在这里强调的核心内容可以概括为：SpecAugment / noise perturbation / augmentation 技术如何改善 ASR 训练。。
从课程组织上看，这一部分承担的是把局部概念嵌回整条识别流水线中的作用，避免读者把公式或模块孤立地记忆。

在本章结构中，这一单元被并入“Augmentation and robustness”，目的是把同一阶段的定义、机制和权衡集中讲清。

\subsection{本章小结}
Augmentation and robustness 这一部分把 Advanced Topics on End-to-End ASR I 中相邻的教学单元串成了一条连续叙事：既交代了这一阶段要解决的问题，也交代了方法、推理或工程权衡为什么要这样组织。

\section{Pronunciation and lexicon issues}
根据课程页主题、公开视频字幕以及本讲的 coverage ledger，Pronunciation and lexicon issues 这一部分主要围绕 pronunciation mapping、lexicon 质量与语言资源不齐时的建模困难。 展开。
写作上保持 coverage-first 原则：先把老师在这一阶段真正讲过的问题设定、模型部件、训练或推理逻辑讲清，再压缩成适合复习的结构。

\subsection{pronunciation mapping、lexicon 质量与语言资源不齐时的建模困难}
从证据位置看，这一部分对应的课程证据区间是 00:13:36,120--00:20:29,760。课程在这里强调的核心内容可以概括为：pronunciation mapping、lexicon 质量与语言资源不齐时的建模困难。。
从课程组织上看，这一部分承担的是把局部概念嵌回整条识别流水线中的作用，避免读者把公式或模块孤立地记忆。

在本章结构中，这一单元被并入“Pronunciation and lexicon issues”，目的是把同一阶段的定义、机制和权衡集中讲清。

\subsection{本章小结}
Pronunciation and lexicon issues 这一部分把 Advanced Topics on End-to-End ASR I 中相邻的教学单元串成了一条连续叙事：既交代了这一阶段要解决的问题，也交代了方法、推理或工程权衡为什么要这样组织。

\section{Multilingual ASR motivation}
根据课程页主题、公开视频字幕以及本讲的 coverage ledger，Multilingual ASR motivation 这一部分主要围绕 multilingual ASR 的动机、shared components 和跨语言共享的边界。 展开。
写作上保持 coverage-first 原则：先把老师在这一阶段真正讲过的问题设定、模型部件、训练或推理逻辑讲清，再压缩成适合复习的结构。

\subsection{multilingual ASR 的动机、shared components 和跨语言共享的边界}
从证据位置看，这一部分对应的课程证据区间是 00:20:25,080--00:27:22,840。课程在这里强调的核心内容可以概括为：multilingual ASR 的动机、shared components 和跨语言共享的边界。。
这一段更像课程的现代延伸：它把传统 ASR 课程内容与 multilingual transfer、speech foundation model、fine-tuning 实践接到了一起。

在本章结构中，这一单元被并入“Multilingual ASR motivation”，目的是把同一阶段的定义、机制和权衡集中讲清。

\subsection{本章小结}
Multilingual ASR motivation 这一部分把 Advanced Topics on End-to-End ASR I 中相邻的教学单元串成了一条连续叙事：既交代了这一阶段要解决的问题，也交代了方法、推理或工程权衡为什么要这样组织。

\section{Benchmarks and related tasks}
根据课程页主题、公开视频字幕以及本讲的 coverage ledger，Benchmarks and related tasks 这一部分主要围绕 benchmarks、评价指标、speech translation 与 direct ID 等多任务关系。 展开。
写作上保持 coverage-first 原则：先把老师在这一阶段真正讲过的问题设定、模型部件、训练或推理逻辑讲清，再压缩成适合复习的结构。

\subsection{benchmarks、评价指标、speech translation 与 direct ID 等多任务关系}
从证据位置看，这一部分对应的课程证据区间是 00:27:20,799--00:33:46,440。课程在这里强调的核心内容可以概括为：benchmarks、评价指标、speech translation 与 direct ID 等多任务关系。。
从课程组织上看，这一部分承担的是把局部概念嵌回整条识别流水线中的作用，避免读者把公式或模块孤立地记忆。

记录系统中的附加说明指出：CS224S supplement anchor: multilingual ASR, Whisper, and foundation-model style evaluation; use this block if you want a modern comparison.

在本章结构中，这一单元被并入“Benchmarks and related tasks”，目的是把同一阶段的定义、机制和权衡集中讲清。
\begin{knowledgebox}{CS224S 补充}
CS224S supplement anchor: multilingual ASR, Whisper, and foundation-model style evaluation; use this block if you want a modern comparison.
\end{knowledgebox}

\subsection{本章小结}
Benchmarks and related tasks 这一部分把 Advanced Topics on End-to-End ASR I 中相邻的教学单元串成了一条连续叙事：既交代了这一阶段要解决的问题，也交代了方法、推理或工程权衡为什么要这样组织。

\section{Whisper and speech foundation models}
根据课程页主题、公开视频字幕以及本讲的 coverage ledger，Whisper and speech foundation models 这一部分主要围绕 Whisper and other speech foundation-model style methods在多语种 / 转写 / 时间戳方面的表现。 展开。
写作上保持 coverage-first 原则：先把老师在这一阶段真正讲过的问题设定、模型部件、训练或推理逻辑讲清，再压缩成适合复习的结构。

\subsection{Whisper and other speech foundation-model style methods在多语种 / 转写 / 时间戳方面的表现}
从证据位置看，这一部分对应的课程证据区间是 00:33:44,480--00:40:32,920。课程在这里强调的核心内容可以概括为：Whisper and other speech foundation-model style methods在多语种 / 转写 / 时间戳方面的表现。。
这一段更像课程的现代延伸：它把传统 ASR 课程内容与 multilingual transfer、speech foundation model、fine-tuning 实践接到了一起。

记录系统中的附加说明指出：CS224S supplement anchor: compare with the Spring 2025 lectures on foundation models and speech recognition beyond English.

在本章结构中，这一单元被并入“Whisper and speech foundation models”，目的是把同一阶段的定义、机制和权衡集中讲清。
\begin{knowledgebox}{CS224S 补充}
CS224S supplement anchor: compare with the Spring 2025 lectures on foundation models and speech recognition beyond English.
\end{knowledgebox}

\subsection{本章小结}
Whisper and speech foundation models 这一部分把 Advanced Topics on End-to-End ASR I 中相邻的教学单元串成了一条连续叙事：既交代了这一阶段要解决的问题，也交代了方法、推理或工程权衡为什么要这样组织。

\section{Fine-tuning and adaptation}
根据课程页主题、公开视频字幕以及本讲的 coverage ledger，Fine-tuning and adaptation 这一部分主要围绕 fine-tuning、adaptation 与整讲总结：如何把大模型迁移到特定 speech 任务。 展开。
写作上保持 coverage-first 原则：先把老师在这一阶段真正讲过的问题设定、模型部件、训练或推理逻辑讲清，再压缩成适合复习的结构。

\subsection{fine-tuning、adaptation 与整讲总结}
从证据位置看，这一部分对应的课程证据区间是 00:40:28,680--00:46:55,359。课程在这里强调的核心内容可以概括为：fine-tuning、adaptation 与整讲总结：如何把大模型迁移到特定 speech 任务。。
这一段更像课程的现代延伸：它把传统 ASR 课程内容与 multilingual transfer、speech foundation model、fine-tuning 实践接到了一起。

记录系统中的附加说明指出：CS224S supplement anchor: fine-tuning base speech models and multilingual adaptation.

在本章结构中，这一单元被并入“Fine-tuning and adaptation”，目的是把同一阶段的定义、机制和权衡集中讲清。
\begin{knowledgebox}{CS224S 补充}
CS224S supplement anchor: fine-tuning base speech models and multilingual adaptation.
\end{knowledgebox}

\subsection{本章小结}
Fine-tuning and adaptation 这一部分把 Advanced Topics on End-to-End ASR I 中相邻的教学单元串成了一条连续叙事：既交代了这一阶段要解决的问题，也交代了方法、推理或工程权衡为什么要这样组织。

\section{总结与延伸}
Advanced Topics on End-to-End ASR I 这一讲在整门《Speech Recognition and Understanding》课程中的位置，是把 Advanced E2E ASR I 讲成一个可复用的建模模块。
从教材化角度看，本讲最重要的不是记住单一术语，而是理解它如何嵌入完整的 automatic speech recognition pipeline：输入语音如何被表示、模型如何被训练、解码如何得到最终转写，以及这些设计分别带来什么假设和权衡。
由于公开源以课程页和公开视频字幕为主、缺少官方 slide PDF，本章在正文里尽量把术语、方法和工程语境解释完整，并把缺失项留在 omission 和 manifest 里，而不是凭空补写。
\end{document}
